\clearpage
\subsection{Grenzwerte trotz Schwingung}
\label{subsec:real_functions_squeeze}
Der Sinus besitzt für $x\to+\infty$ keinen Grenzwert:
Bei $x=\pi/2+2k\pi$ ist sein Wert $1$, bei $x=3\pi/2+2k\pi$ ist er $-1$.
Beide Arten von Eingaben werden beliebig groß. Beschränktheit allein
reicht also nicht aus.

Eine immer kleiner werdende Hülle kann die Schwingung dennoch
auf einen Grenzwert festlegen.
\begin{TheoremBox}{Einschließungsprinzip}{real_functions_squeeze}
  Gilt für die betrachtete Annäherung schließlich
  $u(x)\le f(x)\le v(x)$ und haben beide äußeren Funktionen denselben
  endlichen Grenzwert $L$, dann gilt auch
  \[
    \lim_{x\to a}u(x)=\lim_{x\to a}v(x)=L
    \quad\Longrightarrow\quad
    \lim_{x\to a}f(x)=L.
  \]
  Diese Aussage heißt \emph{Einschließungsprinzip}
  \EnglishTerm{squeeze theorem}.
\end{TheoremBox}
% Handschrift: PDF-Seite 43 (grafische Einschließung an einer endlichen Stelle).
\begin{center}
  \CachedTikzFigure{content/resources/chapters/real_functions/tikz/functions_squeeze_at_point.tex}
\end{center}
Die äußeren Funktionen drücken $f(x)$ auf denselben Grenzwert zusammen.
Stetigkeit der drei Funktionen ist dafür nicht nötig; die Einschließung
und die beiden äußeren Grenzwerte genügen.

% Handschrift: PDF-Seiten 43–46 (beide ausgelassenen Quotienten bei null).
\begin{ExampleBox}{Ein Exponentialquotient mit einer Lücke bei null}{real_functions_exponential_quotient_zero}
  Bestimmen Sie $\lim_{x\to0}(e^x-1)/x$. Direktes Einsetzen liefert
  die unbestimmte Form $0/0$; wir verwenden für $|x|\le1$ die Abschätzung
  \[
    x\le e^x-1\le x+x^2.
  \]
  \begin{center}
    \CachedTikzFigure{content/resources/chapters/real_functions/tikz/functions_exponential_squeeze.tex}
  \end{center}
  \begin{ClozeArea}[34]{9}
    Beim Teilen durch $x$ muss das Vorzeichen berücksichtigt werden:
    \[
      \begin{array}{c@{\qquad}c}
        0<x\le1&-1\le x<0\\[0.6em]
        1\le\dfrac{e^x-1}{x}\le1+x&
        1+x\le\dfrac{e^x-1}{x}\le1.
      \end{array}
    \]
    Alle äußeren Ausdrücke gehen für $x\to0$ gegen $1$. Folglich
    \[
      \lim_{x\to0}\frac{e^x-1}{x}=1.
    \]
    Mit dem ergänzten Wert $f(0)=1$ wird die Lücke stetig geschlossen.
  \end{ClozeArea}
\end{ExampleBox}
\begin{ExampleBox}{Ein Kosinusquotient mit einer Lücke bei null}{real_functions_cosine_quotient_zero}
  Bestimmen Sie $\lim_{x\to0}(\cos x-1)/x$ mit der Zählerabschätzung
  $-x^2\le\cos x-1\le x^2$. Auch hier ist der ursprüngliche Term bei $0$ nicht definiert.
  \begin{center}
    \CachedTikzFigure{content/resources/chapters/real_functions/tikz/functions_cosine_squeeze.tex}
  \end{center}
  \begin{ClozeArea}[34]{7}
    Für beide Vorzeichen von $x\ne0$ können wir gemeinsam schreiben:
    \[
      \left|\frac{\cos x-1}{x}\right|\le |x|
      \quad\Longrightarrow\quad
      -|x|\le\frac{\cos x-1}{x}\le |x|.
    \]
    Wegen $\lim_{x\to0}(-|x|)=\lim_{x\to0}|x|=0$ folgt
    \[
      \lim_{x\to0}\frac{\cos x-1}{x}=0.
    \]
    Der passende Wert für eine stetige Fortsetzung ist $f(0)=0$.
  \end{ClozeArea}
\end{ExampleBox}
Die Zählerabschätzungen werden hier als Hilfsaussagen verwendet;
die Zeichnungen veranschaulichen sie, ersetzen aber keinen Nachweis.
\begin{ExampleBox}{Eine abklingende Schwingung}{real_functions_sinc_infinity}
  Bestimmen Sie das Verhalten von $f(x)=\sin x/x$ für $x\to\pm\infty$.
  \begin{ClozeArea}[34]{6}
    Aus der Beschränktheit des Sinus folgt für $x\ne0$
    \[
      \left|\frac{\sin x}{x}\right|\le\frac1{|x|}
      \quad\Longrightarrow\quad
      -\frac1{|x|}\le\frac{\sin x}{x}\le\frac1{|x|}.
    \]
    Beide Schranken gehen gegen null. Also gilt
    $\lim_{x\to\pm\infty}\sin x/x=0$.
  \end{ClozeArea}
  \begin{center}
    \CachedTikzFigure{content/resources/chapters/real_functions/tikz/functions_squeeze.tex}
  \end{center}
\end{ExampleBox}
\begin{RemarkBox}[remark title=Große Eingaben statt Annäherung an null]
  Hier wurde das Verhalten für große $|x|$ untersucht, nicht für $x\to0$.
  An der Stelle null hilft die Schranke $1/|x|$ nicht.
  Den dortigen Grenzwert behandeln wir bei der Differentialrechnung.
\end{RemarkBox}

% Handschrift: PDF-Seite 47 (Exponentialfunktion gegenüber beliebigen Potenzen).
\begin{ExampleBox}{Exponentielles Wachstum setzt sich gegen jede Potenz durch}{real_functions_exponential_dominance}
  Untersuchen Sie $e^x/x^n$ für $x\to+\infty$ mit festem $n\in\N$.
  Wir verwenden, dass für genügend große positive $x$ die Abschätzung
  $e^x\ge x^{n+1}$ gilt.
  \begin{ClozeArea}[34]{5}
    Durch Division durch $x^n>0$ folgt die untere Schranke
    \[
      \frac{e^x}{x^n}\ge x
      \quad\Longrightarrow\quad
      \lim_{x\to+\infty}\frac{e^x}{x^n}=+\infty.
    \]
    Entsprechend gilt $\lim_{x\to+\infty}x^n/e^x=0$.
  \end{ClozeArea}
  \begin{center}
    \CachedTikzFigure{content/resources/chapters/real_functions/tikz/functions_exponential_power_comparison.tex}
  \end{center}
\end{ExampleBox}
\begin{KeyPointBox}{Exponentialfunktion und Potenz im Grenzvergleich}{real_functions_exponential_power_limits}
  Für jedes feste $n\in\N$ gelten
  \[
    \lim_{x\to+\infty}\frac{e^x}{x^n}=+\infty,\qquad
    \lim_{x\to+\infty}x^n e^{-x}=0,\qquad
    \lim_{x\to-\infty}|x|^n e^x=0.
  \]
  Die letzte Aussage folgt für $t=-x$ aus
  $\lim_{t\to+\infty}t^n e^{-t}=0$.
\end{KeyPointBox}
Später liefert die \emph{Regel von de l'Hospital}
\EnglishTerm{L'Hopital's rule} einen weiteren Zugang zu unbestimmten
Formen wie $0/0$ und $\infty/\infty$.
